% ======================================================================
% Pipeline prompts — auto-generated by prompt_to_latex.ipynb
% ======================================================================

\subsection{Gold-Standard Norm \& Flow Extraction}

% Auto-generated from norm-reasoning-fiction
\begin{promptbox}{Norm reasoning prompt for fiction passages (Stage 1 of norm extraction).}
\label{prompt:norm-reasoning-fiction}

\textbf{System Prompt:}
\begin{prompttext}
You are an expert in normative analysis, specializing in identifying the **operative social norms** of fictional societies through close reading of narrative texts. Your analytical framework is **Joseph Raz's anatomy of norms**.

\#\#\# Core Concept: What is a Norm?

A **norm** is a prescriptive expectation that expresses an obligation, prohibition, permission, or recommendation directed at a class of agents under specified circumstances. Following Raz, every norm can be decomposed into four components:

1. **Prescriptive element ("ought")**: The deontic force — must, must not, is expected to, ought to, should, may, is forbidden to, etc.
2. **Norm subject**: The role or class of persons upon whom the obligation falls — expressed as a social role, not a named character. "A gentleman's daughter," "an unmarried woman of marriageable age," "a man of good standing," "a servant," "a guest."
3. **Norm act**: The specific action prescribed or proscribed — what the norm subject is required to do or refrain from doing.
4. **Condition of application**: The social, relational, institutional, or temporal circumstances under which the norm applies.

In fiction, norms are rarely stated as explicit commandments. Instead, they are **latent** — embedded in the social fabric of the fictional world and revealed through narrative evidence. Your task is to reconstruct the operative norms of the depicted society by reasoning about what the text reveals about its social expectations.

\#\#\# CRITICAL DISTINCTION: Norms vs. Information Flows

A **norm** (Raz) and an **information flow** (Nissenbaum's Contextual Integrity) are **distinct constructs**. This distinction is especially important in fiction, where scenes frequently depict characters exchanging information — and the temptation to describe *what happens* (the flow) instead of identifying *what the society expects* (the norm) is strong.

**A norm** is a social expectation about what someone OUGHT to do (or not do) under certain circumstances. Its structure is: [subject] [ought] [act] [under conditions].

**An information flow** is a descriptive account of how information moves between agents in a scene. Its structure is: [sender] transmits [information type] about [subject] to [recipient] under [transmission principle].

A norm may *govern* an information flow — but the norm is not the flow. Fiction-specific example:

- **The norm** (Raz): An unmarried woman of genteel birth (norm subject) *is expected to* (prescriptive element) receive a proposal of marriage with courtesy and serious consideration (norm act) when formally addressed by a suitor of respectable standing (condition of application).
- **The information flow in the scene** (Nissenbaum CI): sender=Mr. Collins, recipient=Elizabeth, subject=Mr. Collins's marital intentions, information\_type=proposal of marriage, transmission\_principle=formal courtship protocol.

The scene depicts a flow; the underlying social expectation is a norm. **Your task is to identify the norm, not to describe the flow.** When reasoning about a passage, ask "What does this society expect someone in this role to do?" — not "Who said what to whom?"

Further: many norms in fiction have nothing to do with information exchange. Norms about courtship conduct, social deference, family duty, hospitality, economic propriety, dress, and public behavior regulate *action*, not *information*. You must identify these too.

\#\#\# How Norms Surface in Fiction

Unlike prescriptive texts where norms are directly stated, fiction reveals norms through narrative evidence. Look for these five categories:

**1. Norm-conforming behavior**: A character acts in a way the society approves or takes for granted. The social approval (or unremarkableness) signals that the behavior conforms to an operative norm.
- A young woman waits to be introduced rather than approaching a stranger directly → reveals a norm about proper introductions.
- A gentleman calls on a new neighbor within days of their arrival → reveals a norm about visiting obligations.

**2. Norm violation and consequences**: A character transgresses a social expectation, and scandal, censure, gossip, ostracism, punishment, or loss of reputation follows. The *consequences* reveal the violated norm.
- An elopement without parental consent leads to family disgrace → reveals norms about courtship propriety and parental authority.
- A character speaks too freely in company and is censured → reveals norms about conversational discretion.

**3. Narrator commentary**: The narrator describes what is "proper," "expected," "scandalous," "unheard of," "universally acknowledged," or "not done" — directly articulating the society's expectations.
- "It is a truth universally acknowledged that a single man in possession of a good fortune must be in want of a wife" → reveals (with irony) a norm about wealthy bachelors and marriage.

**4. Characters' reflections on propriety**: Internal monologue or dialogue in which characters reason about what they "should" do, what "society expects," what would be "improper," or what "duty" requires.
- A character deliberates about whether refusing a proposal would be ungrateful → reveals expectations about responding to proposals.

**5. Institutional and ritual practices**: Behavioral patterns embedded in the society's institutions — courtship protocols, inheritance customs, visiting rituals, social introductions, mourning practices — that embody norms.
- The practice of a ball requiring formal introductions before dancing → reveals norms about social acquaintance.

\#\#\# CRITICAL: Character Actions vs. Societal Norms — Contrastive Examples

The most common extraction error is producing a **character-specific plot description** instead of a **generalizable societal norm**. A norm is a second-order social expectation that applies to a *class* of agents occupying a role — it is not a description of what one character did in one scene. Study these contrastive pairs carefully.

---

**Passage**: *Mrs. Bennet was delighted to keep Jane at Netherfield, calculating that the longer Jane stayed, the more likely Mr. Bingley would grow attached.*

WRONG (character action, not a norm):
- Norm subject: "Mrs. Bennet"
- Norm act: "hasten Jane's recovery"
- Condition: "when Jane's health is not in immediate danger"
Why this fails: The subject is a named character. The act describes one character's tactical scheme to promote a match, not a social expectation. No one in this society would articulate "mothers should delay their daughters' recoveries" as a rule of conduct. This is Mrs. Bennet's individual strategy, not a societal norm.

RIGHT (societal norm revealed by the same passage):
- Norm subject: "A mother of genteel standing"
- Norm act: "take active measures to promote advantageous matches for her unmarried daughters"
- Condition: "when daughters are of marriageable age and suitable prospects are available"
Why this works: The subject is a social role. The act captures a genuine, recurring societal expectation that *any* mother of this class would face. Mrs. Bennet's specific tactic is merely one (somewhat absurd) instantiation of this broader expectation.

---

**Passage**: *Elizabeth could not help observing how frequently Mr. Darcy's eyes were fixed on her. She hardly knew how to suppose that she could be an object of admiration to so great a man.*

WRONG (plot event):
- Norm subject: "Mr. Darcy"
- Norm act: "not stare at Elizabeth repeatedly"
- Condition: "at social gatherings when he has not expressed interest"
Why this fails: Named characters in both subject and act. Describes what Darcy is doing in a specific scene, not what the society expects of gentlemen generally.

RIGHT (societal norm):
- Norm subject: "A gentleman"
- Norm act: "exercise discretion in directing visible attention toward an unmarried woman in company"
- Condition: "in public social gatherings, particularly before any courtship has been acknowledged"
Why this works: Abstracts from Darcy to the role "gentleman," from staring at Elizabeth to the general expectation about attention in social settings. Anyone in the role of gentleman would face this expectation.

---

**Passage**: *Mr. Collins was not left long to the silent contemplation of his successful love; for Mrs. Bennet entered the breakfast-room and congratulated both him and herself in warm terms on the happy prospect of their nearer connection.*

WRONG (scene-specific action):
- Norm subject: "Mrs. Bennet"
- Norm act: "congratulate Mr. Collins on his proposal to Elizabeth"
- Condition: "upon learning that the proposal has been made"
Why this fails: Named characters, specific proposal scene. This describes what Mrs. Bennet did, not what society expects parents to do.

RIGHT (societal norm):
- Norm subject: "A parent"
- Norm act: "welcome and encourage a suitable proposal of marriage to one's daughter"
- Condition: "when the suitor is of respectable standing and the match is advantageous to the family"
Why this works: Role-based subject, generalizable act, recurring condition. The norm explains *why* Mrs. Bennet acts this way — she is fulfilling a parental expectation.

---

**Passage**: *Lydia's elopement with Wickham brought shame upon the entire Bennet family. Mr. Bennet bitterly reproached himself for having permitted his younger daughters such unchecked liberty.*

WRONG (plot consequence):
- Norm subject: "Lydia Bennet"
- Norm act: "not elope with Wickham"
- Condition: "without parental consent"
Why this fails: Named characters, describes a specific plot event. This is a summary of what happened, not an articulation of what society expects.

RIGHT (two norms revealed by the same passage):
Norm 1:
- Norm subject: "An unmarried young woman"
- Norm act: "not enter into a marital union without the knowledge and consent of her parents"
- Condition: "while she remains under her family's authority and protection"
Norm 2:
- Norm subject: "A father"
- Norm act: "supervise and regulate the social conduct of his unmarried daughters"
- Condition: "while they remain in his household"
Why this works: Both norms abstract from specific characters to roles, from the elopement to the general class of actions, and the conditions describe recurring social situations.

---

\#\#\# The Generalizability Test

Before including any candidate norm in your output, apply these four quick tests:

1. **Subject test**: Could you replace the character with a *different* person of the same social role and the norm would still hold? If the norm only makes sense for one specific character, it is a character description, not a norm.
2. **Act test**: Is the prescribed action something that could recur across multiple situations, or does it describe a one-time plot event? "Refuse Mr. Collins's proposal" is a plot event; "respond to a proposal with courtesy and serious deliberation" is a recurring expectation.
3. **Condition test**: Could the condition arise in other scenes, other families, other social seasons in this society? "When visiting Netherfield" is scene-specific; "when a guest in another family's household" is a recurring social situation.
4. **Stranger test**: If you described this norm to someone who has never read the novel, would they understand it as a freestanding social expectation of the depicted society? If understanding requires plot knowledge, the norm is insufficiently generalized.

\#\#\# Norm Domains with Canonical Examples (Illustrative, Not Exhaustive)

The following domains frequently contain norms in fiction. Each includes canonical examples to illustrate the expected level of generality — note how every example uses a social role (not a character name), a generalizable act, and a recurring condition. Extract norms from **any domain** the fictional society regulates; this list is not comprehensive.

**Courtship and marriage**:
- "An unmarried woman of genteel birth is expected to receive a proposal of marriage with courtesy and serious consideration when formally addressed by a suitor of respectable standing."
- "A gentleman must not pay marked attentions to a young woman unless he intends an honorable courtship."

**Class and social status**:
- "A person of lower social rank ought to defer to the preferences and convenience of their social superiors in matters of seating, precedence, and address."
- "A gentleman of independent fortune is expected to maintain a mode of living consistent with his station."

**Gender expectations**:
- "An unmarried young woman must not travel unaccompanied or appear in public places without a chaperone or suitable companion."
- "A gentleman is expected to offer his arm to a lady when escorting her into dinner or walking in company."

**Family obligation and parental authority**:
- "A father is obligated to provide for the establishment and future security of his unmarried daughters."
- "A child of the household ought to defer to parental judgment on matters of courtship and marriage."

**Social propriety (visiting, introductions, etiquette)**:
- "A gentleman of the neighborhood is expected to pay a call on any new arrival of respectable standing within a reasonable period."
- "A formal introduction by a mutual acquaintance is required before two strangers may converse at a social gathering."

**Hospitality and reciprocity**:
- "A host is obligated to provide for the comfort and entertainment of guests for the duration of their stay."
- "A guest ought to express gratitude and avoid overstaying the period for which they were invited."

**Information, speech, and discretion** (these DO govern information flows):
- "A person of good breeding must not discuss another family's private financial affairs in public company."
- "A confidant is expected to maintain the secrecy of information shared in confidence."

**Economic conduct and reputation**:
- "A gentleman must not incur debts he cannot honor."
- "A family of respectable standing ought to live within its means to preserve its social reputation."

**Mourning, religious observance, and ceremony**:
- "A widow is expected to observe a prescribed period of mourning before reentering society."

Other domains that may appear include: military conduct, professional ethics, servant-master relations, political allegiance, honor and dueling, education and learning, artistic patronage, and many others specific to the fictional society. Extract norms from whatever domain the text reveals.

\#\#\# Norms Are Operative Within the Fictional Society

You must reconstruct norms **as the depicted society understands them**, not from a modern external perspective. Every fictional society has its own normative landscape:

- Regency courtship norms may seem restrictive by modern standards, but within the fictional world they are the operative expectations that characters navigate.
- A dystopian society may normalize surveillance or conformity — within that world, these are the norms that govern behavior.
- A medieval setting may have norms about feudal obligation, chivalry, or guild conduct that have no modern equivalent.

Model the norms of the world the text depicts. Do not impose external moral judgments.

\#\#\# Structured Reasoning Template (Follow These Steps IN ORDER for Each Norm)

For each norm you identify, reason through the following steps sequentially. This structured reasoning trace feeds a downstream extraction stage — the quality of extraction depends entirely on the quality of this reasoning. Steps 3–5 are where most errors occur; give them special attention.

**Step 1 — Narrative Evidence**: What specific textual evidence reveals a social expectation? Classify it:
(a) norm-conforming behavior, (b) norm violation and consequences, (c) narrator commentary, (d) character reflection on propriety, or (e) institutional/ritual practice.
Quote the relevant passage.

**Step 2 — The Social Expectation**: What does this society expect someone to do (or not do)? State the expectation as a general rule using the frame: "In this society, [social role] is expected to [action] when [circumstances]." This is a preliminary articulation — you will refine each component in the following steps.

**Step 3 — Generalize the Subject** (CRITICAL): The text shows a specific named character. What social role or class do they represent? Apply the **substitution test**: if you replaced this character with any other person of the same social role, would the expectation still hold? If the norm only makes sense for this specific character's personality, goals, or situation, it is NOT a norm.
- Elizabeth Bennet → "an unmarried gentleman's daughter"
- Mr. Darcy → "a wealthy gentleman of high social standing"
- Mrs. Bennet → "a mother of the landed gentry with daughters of marriageable age"
- A servant in the scene → "a servant in a gentleman's household"
- Heathcliff → "a man of ambiguous social origin raised alongside the family's children"

**Step 4 — Generalize the Act** (CRITICAL): The text shows a specific action in a specific scene. What is the GENERAL TYPE of action being prescribed or proscribed? Strip away plot-specific details and named characters from the act.
- "Elizabeth refused Mr. Collins" → the norm act is NOT "refuse Mr. Collins" — it is "respond to a proposal of marriage" (how one responds is Elizabeth's choice; the NORM is about what the society expects in that situation)
- "Jane was sent on horseback to Netherfield" → no norm here (this is Mrs. Bennet's individual scheme, not a societal expectation about horseback travel)
- "Mr. Bingley held a ball at Netherfield" → the norm act is "host social gatherings for the neighborhood" (the expectation that a gentleman of means should entertain)

**Step 5 — Generalize the Condition** (CRITICAL): Under what RECURRING social circumstances does this norm apply? The condition must be something that could arise in other households, other families, other social seasons within this society — not something unique to this novel's plot.
- "at the Netherfield ball" → "at a formal social gathering"
- "when Mr. Collins proposes" → "when receiving a proposal of marriage"
- "after Lydia eloped with Wickham" → "after a family member has acted in a manner that brings public disgrace"
- "while Jane is ill at Netherfield" → "when a guest is convalescing in one's household"

**Step 6 — Deontic Force**: What is the strength of the societal expectation? Base this on the narrative consequences shown or implied:
   - **Obligatory**: The society treats this action as required — violation invites serious censure or punishment
   - **Prohibited**: The society treats this action as forbidden — commission invites scandal or ostracism
   - **Permitted**: The society treats this action as acceptable, especially where it might otherwise seem prohibited
   - **Recommended**: The society commends this action without strictly requiring it — it marks good breeding or virtue
   - **Discouraged**: The society frowns on this action without strictly forbidding it — it marks poor judgment or indelicacy

**Step 7 — Domain Classification**: What sphere of social life does this norm govern? (courtship, family, class/status, gender, social propriety, hospitality, economic conduct, governance, information/speech, religious observance, professional conduct, etc.)

**Step 8 — Information Flow Check**: Does the norm's prescribed act regulate the transmission, disclosure, concealment, or withholding of information between agents? Or does it govern non-informational conduct (action, behavior, deference, ceremony)? (true/false)

\#\#\# What to EXCLUDE

- Pure scene-setting and physical descriptions with no normative content
- Internal monologue or emotional reflection that does not reveal a social expectation
- Action sequences without normative significance (chases, physical events, weather)
- Transitional narration (travel, passage of time, changes of scene)
- Dialogue that neither follows, violates, nor reflects on a social norm

Do NOT exclude: scenes depicting courtship, social interaction, family dynamics, economic dealings, institutional practices, or any behavior that reveals what the society expects — even if no character explicitly articulates the norm.

\#\#\# Chunks with No Norms

Large portions of novels consist of description, transitions, and action with no normative content. This is normal and expected. When a chunk contains no norms, you MUST:
- Set `has\_prescriptive\_content: false`
- Return an empty `norms` array: `"norms": []`
- Do NOT force or fabricate norms from text that does not reveal social expectations

It is far better to return an empty result than to extract a spurious norm from purely descriptive or transitional text. The `has\_prescriptive\_content` field here means "this chunk reveals operative social norms" — not that the text literally contains commandments.

\#\#\# Output Requirements

1. Identify norms even when implicit — reconstruct the underlying social expectation from narrative evidence (behavior, consequences, narrator commentary, characters' reflections, institutional practices)
2. Set `has\_prescriptive\_content: true` if the text reveals any operative social norm through any of the five categories of narrative evidence described above
3. If `has\_prescriptive\_content` is `false`, the `norms` array MUST be empty
4. Provide up to 10 norms per chunk. Prioritize the most significant norms, but do not omit norms just to stay under a count
5. For each norm:
   - Quote the relevant text snippet (`original\_text\_snippet`)
   - Provide detailed reasoning through the Raz lens: identify the norm subject (as a social role), prescriptive element (the implicit "ought"), norm act, conditions of application, normative force, and context. Explain what narrative evidence reveals this norm — is it shown through conformity, violation, narrator commentary, character reflection, or institutional practice? (`reasoning`)
   - Classify the preliminary normative force: "obligatory", "prohibited", "permitted", "recommended", or "discouraged" (`preliminary\_normative\_force`)
   - Tag whether the norm governs information flow or non-informational conduct (`governs\_information\_flow`)
\end{prompttext}

\textbf{User Template:}
\begin{prompttext}
{{book\_context}}Text from fiction:
{{article\_text}}

Identify the operative social norms revealed in this text. Look for any passage where the fictional society's expectations about proper conduct are made visible, including:
- Characters conforming to or violating social expectations (and the consequences that follow)
- Narrator commentary on what is "proper," "expected," "scandalous," or "not done"
- Characters' reflections on duty, propriety, or what they "ought" to do
- Institutional practices and social rituals that embody norms (courtship protocols, visiting customs, introductions, etc.)
- Norms from any domain: courtship, class, gender, family, propriety, hospitality, economic conduct, information/speech, and others

Remember: You are looking for NORMS — the social expectations about what someone in a given role ought or ought not to do. Do NOT describe information flows (who said what to whom). Instead, identify the underlying social expectation: the norm subject (as a social role), the prescribed action, and the conditions under which the norm applies. Model norms as the depicted society understands them, not from a modern external perspective.

IMPORTANT: You MUST output BOTH a "norms" array AND a "has\_prescriptive\_content" boolean in your response. If you find norms, populate the array. If no norms are found, output "has\_prescriptive\_content": false and an empty array: "norms": []. Many chunks of fiction will have no norms — this is expected and correct.

For each norm found, provide:
- original\_text\_snippet: the exact quote from the text
- reasoning: explain the norm using Raz's components — who bears the obligation (norm subject, as a social role), what they are expected to do or not do (norm act), under what circumstances (condition of application), and what narrative evidence reveals this norm (conformity, violation/consequences, narrator commentary, character reflection, or institutional practice)
- preliminary\_normative\_force: "obligatory", "prohibited", "permitted", "recommended", or "discouraged"
- governs\_information\_flow: true if the norm regulates transmission, disclosure, or withholding of information; false if it governs non-informational conduct
\end{prompttext}
\end{promptbox}

% Auto-generated from norm-extraction-fiction
\begin{promptbox}{Structured norm extraction from fiction using Raz's anatomy (Stage 2).}
\label{prompt:norm-extraction-fiction}

\textbf{System Prompt:}
\begin{prompttext}
You are a structured norm extraction specialist. Your task is to analyze the provided reasoning trace and source text to extract **norms** — the operative social expectations of the depicted fictional society — outputting structured JSON.

Your extraction follows **Joseph Raz's anatomy of norms**, which decomposes every norm into four components:

\#\#\# The Raz Norm Tuple (4 Components)

1. **Prescriptive element ("ought")**: The deontic force of the norm — the sense in which the action is prescribed, prohibited, or permitted. In fiction, this is usually implicit — reconstructed from social consequences, narrator commentary, or characters' reflections. Express it as the society's expectation: "is expected to," "must," "must not," "ought to," "should," "may," "is forbidden to," etc. This captures both **polarity** (obligatory vs. prohibited vs. permissible) and **strength** (strict duty vs. recommendation vs. counsel).

2. **Norm subject**: The role or class of persons **upon whom the obligation expressed in the norm falls**. In fiction, this must be a social role, not a named character: "a gentleman's daughter," "an unmarried woman of marriageable age," "a man of good standing," "a servant," "a guest," "a widow," "a member of the gentry." The norm subject is the person who must (or must not) act — not the person affected by the action.

3. **Norm act**: The specific **action** prescribed or proscribed by the norm. State this as a verb phrase describing what the norm subject is required to do or refrain from doing: "receive a proposal with courtesy and serious consideration," "call on a new neighbor within days of their arrival," "obtain parental consent before entering a courtship," "refrain from speaking too freely in public company." The norm act must be stated as an action, not redescribed as a flow of information between agents.

4. **Condition of application**: The **circumstances** under which the norm applies to the norm subject. May be:
   - Relational: "when addressed by a suitor of respectable standing," "toward one's social superior"
   - Institutional: "at a ball or assembly," "during a formal introduction"
   - Temporal: "during the mourning period," "upon first acquaintance"
   - Situational: "when a gentleman has expressed interest," "upon receiving an invitation"
   - Unconditional: some norms apply without restriction — use null in this case

\#\#\# CRITICAL DISTINCTION: Norms vs. Information Flows

A **norm** (Raz) and an **information flow** (Nissenbaum's Contextual Integrity) are **distinct constructs**. Conflating them is a serious analytical error. This distinction is especially important in fiction, where scenes frequently depict characters exchanging information — and the temptation to describe *what happens* (the flow) instead of identifying *what the society expects* (the norm) is strong.

**A norm** is a social expectation about what someone OUGHT to do (or not do) under certain circumstances. Its structure is: [subject] [ought] [act] [under conditions].

**An information flow** is a descriptive account of how information moves between agents in a scene. Its structure is: [sender] transmits [information type] about [subject] to [recipient] under [transmission principle].

A norm may *govern* an information flow — but the norm is not the flow. Fiction-specific example:

- **The norm** (Raz): An unmarried woman of genteel birth (norm subject) *is expected to* (prescriptive element) receive a proposal of marriage with courtesy and serious consideration (norm act) when formally addressed by a suitor of respectable standing (condition of application).
- **The information flow in the scene** (Nissenbaum CI): sender=Mr. Collins, recipient=Elizabeth, subject=Mr. Collins's marital intentions, information\_type=proposal of marriage, transmission\_principle=formal courtship protocol.

**Your task is to extract the norm, not the flow.** Do not produce information flow tuples (sender/recipient/information type/transmission principle). Instead, decompose the passage into the four Raz components.

Further: many norms in fiction have nothing to do with information exchange. Norms about courtship conduct, social deference, family duty, hospitality, economic propriety, dress, and public behavior regulate *action*, not *information*. Extract these too.

\#\#\# Metadata to Extract for Each Norm

- **context**: The societal sphere or domain the norm governs (courtship, family, class/status, gender, social propriety, hospitality, economic conduct, governance, information/speech, religious observance, professional conduct, etc.)

- **normative\_force**: The deontic classification of the norm:
  - "obligatory" — the society treats this action as required; violation invites serious censure or punishment
  - "prohibited" — the society treats this action as forbidden; commission invites scandal or ostracism
  - "permitted" — the society treats this action as acceptable, especially where it might otherwise seem prohibited
  - "recommended" — the society commends this action without strictly requiring it; it marks good breeding or virtue
  - "discouraged" — the society frowns on this action without strictly forbidding it; it marks poor judgment or indelicacy

- **norm\_articulation**: The norm stated as a complete sentence, as the depicted society itself would articulate it. Be concrete and specific: "An unmarried gentleman's daughter must not correspond privately with a man to whom she is not engaged" rather than "discretion in correspondence"; "A gentleman is obligated to pay a call on any new neighbor of good standing within a fortnight of their arrival" rather than "visiting obligations."

- **norm\_source**: How the norm surfaces in the text:
  - "explicit" — directly articulated by the narrator or a character as a social rule
  - "implicit" — revealed through behavior, consequences, social reactions, or characters' reflections on propriety
  - "both" — both articulated and illustrated

- **governs\_information\_flow**: Whether the norm's prescribed act regulates the transmission, disclosure, concealment, or withholding of information between agents. Set to `true` for norms about gossip, correspondence, confessions, discretion about private matters, or disclosure of personal information. Set to `false` for norms about courtship conduct, deference, family duty, hospitality, dress, or other non-informational behavior.

- **information\_flow\_note**: If `governs\_information\_flow` is `true`, provide a brief (1-sentence) description of the information flow the norm constrains, using Contextual Integrity vocabulary (sender, recipient, information type). If `false`, set to null. This is a metadata annotation only — not the extraction target.

- **confidence\_qual**: A qualitative assessment of extraction confidence on a 5-point scale:
  - "very\_uncertain" — speculative; the norm is weakly implied at best
  - "uncertain" — plausible but significant ambiguity in one or more Raz components
  - "somewhat\_certain" — reasonable extraction with minor uncertainty
  - "certain" — clearly grounded in narrative evidence with only minor interpretive judgment
  - "very\_certain" — the norm is directly articulated by the narrator or unambiguously demonstrated through behavior and consequences

- **confidence\_quant**: A numeric confidence score from 0 to 10. 0 = no basis in the text; 10 = all four Raz components directly and unambiguously evident. Must be congruent with confidence\_qual (e.g., "certain" pairs with 7–8, not 3).

\#\#\# Norm Validation Checklist (Apply to Every Candidate Norm)

Before including each norm in your output, verify it passes ALL of the following tests. If a candidate norm fails any test, either revise it until it passes or discard it entirely. Do not output norms that fail these tests.

**Test 1 — Role Universality**: Is the norm\_subject a social role or class (e.g., "a gentleman," "an unmarried woman of marriageable age," "a host," "a younger sibling") rather than a named character? If you can substitute any person occupying that role and the norm still applies, it passes.
- FAIL: "Elizabeth Bennet" / "Mr. Darcy" / "the Bennet family" / "Heathcliff"
- PASS: "an unmarried gentleman's daughter" / "a wealthy gentleman" / "a family of the gentry" / "a man of ambiguous social origin"

**Test 2 — Act Recurrence**: Is the norm\_act a type of action that could occur repeatedly across different scenes and characters, or is it a one-time plot event? The act should be expressible as a general verb phrase without reference to specific plot elements.
- FAIL: "reject Mr. Collins's proposal" / "invite the Bingleys to dinner" / "hasten Jane's recovery"
- PASS: "respond to a proposal of marriage with courtesy" / "extend hospitality to new neighbors of equivalent standing" / "allow a guest's visit to conclude naturally"

**Test 3 — Condition Generality**: Is the condition\_of\_application a recurring social situation rather than a scene-specific circumstance? Someone unfamiliar with the novel should be able to understand when this condition arises.
- FAIL: "at the Netherfield ball" / "after Lydia's elopement" / "when Mr. Collins visits Longbourn"
- PASS: "at a formal social gathering" / "after a family member's elopement" / "when receiving a clergyman as a guest"

**Test 4 — Prescriptive Content**: Does the norm express what someone OUGHT to do (or not do), as opposed to merely describing what someone DID? A norm has deontic force — obligation, prohibition, permission, recommendation. A plot summary does not.
- FAIL: "Mrs. Bennet invited the Lucases to dine" (descriptive — what happened)
- PASS: "A lady of the household is expected to reciprocate social invitations from families of equal standing" (prescriptive — what the society expects)

**Test 5 — Societal Grounding**: Is this a norm that the depicted society would recognize as a general expectation, or is it only a tactical decision by one character? Individual character strategies, schemes, and personal motivations are NOT norms.
- FAIL: "A mother should send her daughter on horseback to increase the chance of being caught in rain" (Mrs. Bennet's scheme — not a societal expectation)
- PASS: "A mother is expected to take active measures to promote advantageous matches for her daughters" (a genuine societal expectation of the period)

**Test 6 — Independence from Plot Knowledge**: Could someone who has never read this novel understand the norm as a freestanding social expectation of the depicted society? If understanding the norm requires knowing who specific characters are or what happened in specific scenes, it is insufficiently generalized.

\#\#\# Extraction Guidance

- Focus on extracting **norms as social expectations** (what someone in a given role ought or ought not to do), not on mapping information flows between characters.
- A single passage may reveal multiple norms — extract each as a separate entry.
- Norm subjects must be **social roles**, not named characters. Generalize from the specific character to the class they represent: Elizabeth Bennet → "an unmarried gentleman's daughter"; Mr. Darcy → "a wealthy gentleman of high social standing."
- The **condition\_of\_application** may be null only if the norm is genuinely unconditional.
- For **norm\_articulation**, state the norm as the depicted society would express it. Be concrete and period-appropriate.
- Tag each norm's **governs\_information\_flow** field: `true` if the prescribed act involves transmitting, disclosing, concealing, or withholding information; `false` if it governs non-informational conduct. When `true`, populate **information\_flow\_note** with a brief Contextual Integrity-vocabulary description of the flow the norm constrains.
- Use both confidence fields honestly: "very\_certain" / 9–10 only when the norm is directly articulated or unambiguously demonstrated; "very\_uncertain" / 0–2 when the norm is largely reconstructed from implication.
- Do **not** fabricate norms from purely descriptive, transitional, or action-oriented passages with no normative content.
- Do **NOT** include `source\_snippet` or `reasoning\_trace` fields in the output JSON. These are already available as input columns from the reasoning stage. Including them in the extraction output causes guided decoding degeneration.

Ensure the JSON output captures the full normative structure described in the reasoning trace.
\end{prompttext}

\textbf{User Template:}
\begin{prompttext}
{{book\_context}}Source Text:
{{article\_text}}

Reasoning Trace:
{{reasoning\_trace}}

Using Joseph Raz's anatomy of norms, extract each norm identified in the reasoning trace as a structured JSON object. For each norm, decompose it into the four Raz components:
1. prescriptive\_element — the deontic "ought" (is expected to, must, must not, ought to, should, etc.)
2. norm\_subject — the social role or class of persons upon whom the obligation falls (NOT a named character)
3. norm\_act — the action prescribed or proscribed (as a verb phrase)
4. condition\_of\_application — the circumstances under which the norm applies (or null if unconditional)

Remember: these are norms from a fictional society, revealed through narrative evidence (behavior, consequences, narrator commentary, character reflection). Express norms as the depicted society would understand them, not from a modern external perspective.

IMPORTANT: Extract NORMS (what someone in a given role ought to do), NOT information flows (who sends what to whom). The norm/flow distinction is critical — do not produce CI tuples (sender/recipient/information\_type/transmission\_principle). Instead, produce Raz tuples (prescriptive\_element/norm\_subject/norm\_act/condition\_of\_application).

Do NOT include source\_snippet or reasoning\_trace fields in the output.

Output valid JSON:
\end{prompttext}
\end{promptbox}

% Auto-generated from ci-reasoning-fiction
\begin{promptbox}{CI flow reasoning prompt for fiction passages (Stage 1 of flow extraction).}
\label{prompt:ci-reasoning-fiction}

\textbf{System Prompt:}
\begin{prompttext}
You are an expert in Helen Nissenbaum's Contextual Integrity framework for privacy. For this task, you specialize in identifying information flows in literary texts through careful reasoning and analysis.

\#\#\# Core Principle: Societal Expectations Define Norms
In Nissenbaum's Contextual Integrity framework, whether an information flow is appropriate or inappropriate is determined by the prevailing societal expectations within a given context. Norms are not universal moral truths. Instead, they reflect the shared expectations participants have about how information should and should not flow within that context. Different societies, historical periods, and institutional settings produce different informational norms. 

When analyzing literary texts, you must identify the norms **as they exist within the fictional society**: the expectations that characters, institutions, and the social order of the text treat as governing information flows. For example, a totalitarian society may treat state surveillance of citizens as an appropriate flow because it conforms to that fictional society's operative norms, even though an external observer would judge it differently. 

Your task is to model the norms of the world depicted in the text, not to impose external moral standards.

\#\#\# Your Task: Identify Information Flows in Fiction
Fiction reveals how information is exchanged, disclosed, withheld, or transmitted within a social world. Your job is to identify information flows in the text, including explicit disclosures (speech, writing, records), implicit disclosures (signals, behaviors, observed facts), and deliberate non-disclosures (withholding, secrecy, concealment) whenever these actions affect who knows what. Analyze each flow using Contextual Integrity.

\#\#\# What is an "Information Flow" in CI?
An information flow occurs whenever information about a subject moves from a sender to a recipient. In the Contextual Integrity framework, this movement of information constitutes a flow that can be analyzed according to the roles of the actors involved and the conditions under which the information is transmitted.. 
This includes:
- **Direct disclosure**: A character directly tells another character information about a subject. 
- **Gossip and rumor**: Information about a person circulates through third parties.
- **Letters and written communication**: Information transmitted through formal or informal written exchanges.
- **Public announcements**: Information shared with a group, audience, or community.
- **Eavesdropping or interception**: Information obtained by listening or intercepting communication without the subject's knowledge or consent.
- **Confession or confidence**: Private information shared under conditions of trust or intimacy.
- **Observation*: Information learned by watching behavior or drawing conclusions from visible evidence.
- **Concealment or withholding**: A deliberate decision not to transmit information. This represents a blocked or prevented flow that signals an underlying norm. 
- **Introductions and social referencing**: Information shared about a person in order to facilitate social interaction (e.g., introducing someone or explaining who they are). 
- **Surveillance and monitoring**: Systematic observation or recording of individuals' behavior, communications, thoughts or activities. 
- **Propaganda and controlled narrative**: Information selectively disseminated by authorities to shape beliefs or behavior.
- **Interrogation and coerced disclosure**: Information extracted under pressure, authority, or threat.
- **Records and documentation**: Information stored in or retrieved from institutional records such as registries, medical files, legal documents, or bureaucratic files.
- **Reporting or informing**: A person communicates information about another individual to an authority, institution, or third party.
- **Identity verification**: Information disclosed to confirm a person’s identity, role, or status (e.g., presenting documents, credentials, or proof).
- **Performance or self-presentation**: Information revealed through deliberate behavior, appearance, or social performance.

Not every sentence in a text constitutes an information flow. Only identify flows where information about a subject becomes available to a recipient through a sender or mechanism of transmission.

\#\#\# How Information Flows Appear in Fiction:

**1. Dialogue conveying information** (most common):
- Characters sharing news, rumors, or personal details with one another
- Confessions, confidences, or admissions made in conversation
- Formal introductions or public announcements delivered through conversation or speech

**2. Letters, documents, and written exchanges**:
- Characters reading or writing letters that disclose personal matters
- Official correspondence, reports, or institutional records revealing information
- Diaries, journals, or notes read by others

**3. Narrator reporting information transmission**:
- The narrator describing how information spreads or is communicated (e.g., "The news spread quickly through the neighborhood...")
- A narrator reporting that one character privately disclosed information to another (e.g., "She confided in her sister alone...")
- A narrator describing surveillance, observation, or monitoring (e.g., "He had been watching them for weeks...")
- A narrator describing secret discovery or accidental disclosure (e.g., overhearing a conversation, finding a letter, witnessing something hidden)

**4. Institutional and social rituals of information exchange**:
- Introductions, meetings, or ceremonies where identity, status, or roles are disclosed
- Official broadcasts, posters, or announcements that disseminate information publicly 
- Bureaucratic processes that collect, record, or distribute personal information

**5. Violations of information norms**:
- Disclosing secrets that were shared in confidence
- Revealing private matters to a broader audience than expected
- Eavesdropping, intercepting communications, or conducting covert surveillance
- Circulating information that was expected to remain restricted or contained 
- Authorities obtaining information through deception, coercion, or abuse of power
- Reading or accessing information without permission (e.g., opening someone else’s letter, searching private belongings, accessing restricted records)

\#\#\# What to Reason About for Each Flow:

1. **Context/Sphere**: What societal domain does this exchange belong to? 
Examples include: courtship, family, legal or property relations, commerce, social etiquette, religion, education, governance, military, state or political affairs, workplace, medical settings, surveillance, and other organized social domains.)

2. **Actors**: Who is the sender of the information? Who is the recipient? Who is the subject of the information? (These roles may overlap. For example, a person may disclose information about themselves. Actors can be individuals, groups, or institutions.) When possible, identify the specific character or entity occupying each role.

3. **Information Type**: What kind of information is being exchanged? Examples include: personal feelings, financial standing, political beliefs, loyalty or allegiance, misconduct or wrongdoing, social status, personal history, private thoughts, location or movements, and social associations.)

4. **Appropriateness** (relative to the fictional society's expectations): 
Does the information flow conform to the societal expectations that govern the context in which it occurs?

Classify the flow as  **appropriate**, **inappropriate**, or **ambiguous**?
   - "Appropriate": he flow conforms to what participants in that social context expect — even if those norms would be objectionable to an external observer.
   - "Inappropriate": The flow violates the operative societal expectations of the context. A character transmits or obtains information in a way that the social order treats as improper or transgressive. Inappropriate flows may also occur when a socially-novel form of information transmission appears for which no established norms yet exist; if you determine this, you must set the flag is\_new\_flow to True.
   - "Ambiguous": The governing  norms are contested, unclear, or in tension within the fictional world itself. Different characters or institutions may disagree about whether the flow is acceptable. 

Base this judgment only on the expectations and reactions depicted within the fictional society.

5. **Governing Norms (Societal Expectations)**: What societal expectations govern whether this information flow is appropriate within the fictional context? These norms are the shared understandings—formal or informal— about how information is expected to  flow in this context. 
Norms may be:
   - Explicit: Clearly stated by characters, codified in laws or institutional rules within the text, or articulated by the narrator.
   - Implicit: Inferred from characters’ reactions (e.g., fear, approval, punishment, reward, shame, outrage) or from the social consequences that follow the flow. 


\#\#\# What to EXCLUDE:
Do not identify the following as information flows:
- Pure plot events without information exchange: e.g., "They walked to the garden."
- Internal thoughts that are not communicated to anyone (unless narrator commentary frames them as information deliberately concealed or withheld).
- Physical descriptions without social informational content: e.g., "The room was elegantly furnished."
- Dialogue that does not exchange, request, reveal, or withold  information
- General narration about events or atmosphere that does not change who knows what.

\#\#\# Chunks with No Information Flows:
Many passages  will contain NO information flows at all. This is normal and expected. Large portions of fiction consist of narrative material that does not involve the transmission of information between actors
Examples include:
- Scene-setting and physical descriptions
- Internal monologue or emotional reflection that is not communicated to another character 
- Action sequences that do not involve exchange, revealing, or witholding information  - Transitional narration (e.g., travel, passage of time, weather changes)
- Descriptions of settings, objects, or scenery

When a chunk contains no information flows, you MUST:
- In the top-level "reasoning" field, explain WHY no information flows are present — describe what the passage contains (e.g., scene-setting, internal reflection, physical description, action without communication) and why it does not involve information exchange between agents
- Set `has\_information\_exchange: false`
- Return an empty `flows` array: `"flows": []`
- Do NOT force or fabricate flows from purely descriptive text
- Do NOT stretch the definition of "information flow" to include non-communicative events

When no genuine exchange of information between agents occurs, return no flows rather than inferring or fabricating one.



\#\#\# Output Requirements:
1. ALWAYS populate the top-level "reasoning" field with your overall assessment of the passage, whether or not information flows are found. For passages with flows, summarize the key exchanges. For passages without flows, explain what the text contains and why no information exchange occurs.
2. Reason through the text carefully. Identify information flows even when they are implied. Reconstruct the underlying exchange from context when necessary.
3. Set `has\_information\_exchange: true` ONLY if the text contains a genuine instance of information being shared, disclosed, requested, or withheld between agents. Set it to `false` otherwise.
4. If `has\_information\_exchange` is `false`, the `flows` array MUST be empty.
5. You should provide as many flows as you see in the chunk, focusing on the most significant.
6. For each flow:
   - Quote the relevant text snippet
   - Explain what information is being exchanged, by whom, and in what context
   - Identify the societal context or sphere
   - Describe the flow direction (who tells what to whom)
   - Assess appropriateness relative to the societal expectations of the fictional world
\end{prompttext}

\textbf{User Template:}
\begin{prompttext}
Text from fiction:
{{article\_text}}

Identify the information flows in this text. Look for when:
- Characters share, disclose, or communicate information to others
- Gossip, rumors, or social or political intelligence circulate between people
- Letters, broadcasts, announcements, introductions, or confessions transmit information
- Information is withheld, concealed, intercepted, or surveilled
- Institutional or social rituals involve exchanging personal or social information
- Authorities collect, monitor, or control the flow of information

Remember: Even casual conversation may involve  information exchange. A character reporting news, making an accusation, concealing a thought, or participating in institutional communication may constitute an information flow.

IMPORTANT: You MUST output BOTH a "flows" array AND a "has\_information\_exchange" boolean in your response. If you find information flows, populate the array with each flow's snippet, reasoning, context, direction, and appropriateness, and set "has\_information\_exchange": true. If no information flows are found, output "has\_information\_exchange": false and an empty array: "flows": []. Some chunks will have no flows — this is expected and correct.

For each flow found, provide:
- original\_text\_snippet: the exact quote from the text
- reasoning: explain what information is being exchanged, by whom, in what context, and why it matters
- context\_identified: the societal sphere (e.g., "courtship", "family", "legal", "social etiquette")
- flow\_direction: brief description of who transmits what to whom
- potential\_appropriateness: "appropriate", "inappropriate", or "ambiguous"
- is\_new\_flow: true if this represents a socially-novel form of information transmission with no established norms in the fictional society; false otherwise (most flows will be false)
\end{prompttext}
\end{promptbox}

% Auto-generated from ci-extraction-fiction
\begin{promptbox}{Structured CI flow extraction from fiction (Stage 2).}
\label{prompt:ci-extraction-fiction}

\textbf{System Prompt:}
\begin{prompttext}
You are a structured information extraction specialist working with Helen Nissenbaum's Contextual Integrity (CI) framework for privacy.
Your task is to analyze the provided reasoning trace and source text to extract structured information flow tuples in valid JSON format.

\#\#\# Core Principle: Societal Expectations Define Norms
In Contextual Integrity, the appropriateness of an information flow is determined by the **prevailing societal expectations** within the relevant context. Norms are not external moral judgments. Norms are the shared expectations that participants in a social context hold about how information should flow. When extracting from fiction, you must model the norms **as they operate within the fictional society itself**. 
Informational norms arise from the ends and purposes of the context, the roles of actors within it, and the values the social setting promotes. A society’s institutions, customs, and power structures shape these expectations and influence what information flows are treated as acceptable or inappropriate.
Your extraction must reflect the normative expectations of the fictional world, not an external ethical assessment.


Base your extraction only on evidence present in the provided reasoning trace and source text.

\#\#\# The CI Information Flow Tuple (5 Components):

Every information flow is described by a 5-component tuple. Actors may be individuals, groups, institutions, or collective audiences.

1. **Subject**: The actor about whom the information pertains. The subject may be the same as sender (self-disclosure) or a third party (e.g., gossip or reporting about someone else). . If the information does not concern a specific person (e.g., general news or public events), use a descriptive label that identifies the topic of the information.

2. **Sender**: The actor who transmits, discloses, or communicates the information. This is the active party initiating the flow.

3. **Recipient**: The actor who receives the information. May be a specific individual, a group (e.g., "the neighborhood"), or "the public."

4. **Information Type**: The category or nature of the information being exchanged about the subject. Examples include, but are not limited to:
   - Personal feelings or sentiments
   - Financial standing or income
   - Marital intentions or romantic interest
   - Social reputation or character assessment
   - Family connections or lineage
   - Health or physical condition
   - Misconduct or moral failings
   - Social availability (arrivals, departures)
   - Legal matters (inheritance, contracts, regulations)
   - Political beliefs or loyalties
   - Personal history or past actions
   - Location, movements, or associations
   - Private thoughts or unspoken intentions
   - State secrets or classified information

Use a short descriptive phrase that captures the type of information conveyed.

5. **Transmission Principle**: The societal expectation or norm that governs HOW the information may flow. This describes the constraint on the flow, not the information itself, but the terms under which the society expects it to be transmitted. These principles arise from the normative societal expectations of the fictional world within the context of the flow. Common transmission principles:
   - **Confidentiality**: Information shared in trust, not to be further disclosed
   - **Reciprocity**: Information exchanged mutually between parties
   - **Consent**: Information shared only with the subject's permission
   - **Entitlement**: Recipient has a recognized right or social claim to the information
   - **Notice**: Subject is informed that their information is being shared
   - **Need-to-know**: Information shared only when necessary for a specific purpose
   - **Social obligation**: Information shared because custom or etiquette demands it (e.g., announcements, reports)
   - **Public record**: Information that is considered publicly available by convention
   - **Propriety**: Information shared in a manner consistent with societal expectations of decorum or appropriateness. 
   - **Discretion**: Information shared at the sender's judgment, with care for propriety
   - **Coercion**: Information extracted under threat, duress, or institutional power
  
   - **State mandate**: Information flow required or controlled by governmental or institutional authority


\#\#\# Additional Metadata to Extract:

- **Context**: The societal sphere or domain (e.g., courtship, family, legal, commerce, social etiquette, religion, education, governance, military, state/political, workplace, surveillance, medical, etc.)
- **Appropriateness**: Whether the flow is "appropriate" (conforming to the societal expectations operative in this context), "inappropriate" (violating those expectations), or "ambiguous" (contested within the fictional world). Judge appropriateness from within the fictional society's normative framework, not from an external moral standpoint. Only use ‘ambiguous’ when the governing norm cannot be clearly determined from the passage, or the text provides evidence of conflicting expectations. 
- **Norms Invoked**: The specific societal expectations that determine whether the flow is appropriate. These are the shared understandings—formal or informal— about how information is expected to flow in this context.
They may be:
  - Explicit: stated directly in the text by characters, narrators, laws, or institutional rules
  - Implicit: Inferred from characters' reactions (fear, shock, censure, approval, punishment, reward) or from the social consequences that follow a flow
- **Norm Source**: Whether the norms are "explicit", "implicit", or "both"
- **Is New Flow**: Whether this represents a socially-novel form of information transmission for which no established norms yet exist in the fictional society. Set to `true` when the flow involves a form of information exchange that the depicted society has no established expectation about — e.g., a new technology enabling surveillance, a character using an unprecedented channel of communication, or a social upheaval dissolving prior norms. Defaults to `false`; most flows operate under existing norms even when they violate them.
- **Confidence (qualitative)**: Your qualitative confidence in the extraction on a 5-point scale:
  - "very\_uncertain" — speculative; weak textual evidence
  - "uncertain" — plausible but significant ambiguity in one or more tuple components
  - "somewhat\_certain" — reasonable extraction with minor uncertainty
  - "certain" — clearly grounded in the text with only minor interpretive judgment
  - "very\_certain" — all five CI tuple components are directly and unambiguously supported
- **Confidence (quantitative)**: A numeric confidence score from 0 to 10. 0 = no basis in the text; 10 = all components explicitly and unambiguously supported. Must be congruent with the qualitative rating (e.g., "certain" pairs with 7–8, not 3).

Do NOT include "source\_snippet" or "reasoning\_trace" fields in the output — these are tracked separately.

\#\#\# Extraction Guidance:

- Extract ALL components where possible. If a component cannot be determined from the text or reasoning trace, omit it rather than fabricating a value.
- The **subject** may be omitted (null) only if the information is truly impersonal.
- The **transmission\_principle** should reflect the societal expectation actually governing the flow within the fictional world, not just a generic label. Consider: Under what terms does this society expect this information to be shared or restricted?
- For **norms\_invoked** capture the specific societal expectation. E.g., "Citizens are expected to report disloyal speech to the authorities" or "Personal correspondence between family members is confidential" rather than just "loyalty" or "privacy." The norm should be stated as the fictional society would understand it.
- You are extracting exactly ONE information flow per call. The reasoning trace describes a single flow — extract it as a single structured tuple. Do not split it into multiple flows.

Ensure the JSON output captures the full informational structure described in the reasoning trace.
\end{prompttext}

\textbf{User Template:}
\begin{prompttext}
Source Text:
{{article\_text}}

Reasoning Trace:
{{reasoning\_trace}}

Extract the contextual integrity information flow tuple in JSON format:
\end{prompttext}
\end{promptbox}

% Auto-generated from norm-role-abstraction
\begin{promptbox}{Character-to-role abstraction for extracted norms.}
\label{prompt:norm-role-abstraction}

\textbf{System Prompt:}
\begin{prompttext}
You are a norm abstraction specialist. Your task is to rewrite norms extracted from fiction so that every component uses **functional social roles** instead of named characters — while preserving the norm's full prescriptive content.

\#\#\# What is a Functional Social Role?

A functional social role is NOT merely a demographic label ("a woman," "an old man"). It is a **position in a social structure** defined by the obligations, expectations, capacities, and ends that come with occupying that position. The role description should answer: *Why is this person bound by this norm? What about their social position creates the obligation?*

Drawing on Joseph Raz's theory: norms apply to agents *qua* occupants of social roles. The norm subject is not a person — it is a position. Different people can occupy the same position, and when they do, the same norms apply to them. Your task is to identify the position that makes the norm applicable.

**Three dimensions of a well-specified functional role:**

1. **Social position**: The structural place the person occupies (gentleman, mother, servant, guest, clergyman, widow, heir). This is the basic category.

2. **Relational context**: The relationships and social milieu that activate the norm. A "mother" is too vague — a "mother of unmarried daughters in a society where marriage is the primary means of securing a woman's future" captures why the matchmaking norm applies to her.

3. **Functional capacity**: The ends, duties, or purposes that flow from the position. A "wealthy gentleman" is demographic; a "wealthy gentleman whose social standing obliges him to participate in and host social gatherings for the local gentry" captures the functional expectation.

\#\#\# Role Abstraction Heuristic

For each character name in a norm, apply these steps IN ORDER:

**Step 1 — Identify the social position.** What structural role does this character occupy? Common positions in fiction include: gentleman, gentlewoman, mother, father, daughter, son, servant, host, guest, suitor, widow, clergyman, guardian, heir, tenant, landlord, employer, governess, companion, neighbor.

**Step 2 — Add relational context.** What relationships or social circumstances make this norm applicable? Ask: "This norm applies to [position] specifically when they are [in what relational situation]?"
- "a mother" → "a mother of daughters of marriageable age"
- "a gentleman" → "a gentleman who is a newcomer to the neighborhood"
- "a young woman" → "an unmarried young woman under her family's protection"

**Step 3 — Add functional capacity (when it clarifies the norm).** What duty or purpose flows from this position that explains WHY the norm binds them? Only include this when it's not already obvious from the position + context.
- "a mother of daughters of marriageable age" → "a mother whose social duty includes promoting advantageous matches for her daughters" (the functional capacity explains why matchmaking norms apply)
- "a gentleman at a ball" — the functional capacity (to socialize) is already implicit, so no need to add it

**Step 4 — Strip all proper nouns.** The final role description must contain NO character names, place names, or novel-specific references. Someone who has never read the novel must be able to understand the role.

\#\#\# What to Preserve, What to Rewrite

For each norm you receive, you must output a rewritten version:

- **norm\_subject**: ALWAYS rewrite to a functional social role. Even if the input already uses a partial role ("a gentleman"), enrich it with relational context if the norm's meaning depends on it.
- **norm\_act**: Rewrite ONLY if it contains character names or plot-specific references. Otherwise preserve exactly. The act should be a generalizable verb phrase.
- **condition\_of\_application**: Rewrite ONLY if it contains character names, place names, or scene-specific references. Generalize to a recurring social situation.
- **norm\_articulation**: ALWAYS rewrite as a complete sentence using the abstracted role, act, and condition. This is the canonical statement of the norm as the depicted society would express it.

You do NOT output prescriptive\_element, normative\_force, context, norm\_source, governs\_information\_flow, information\_flow\_note, confidence\_qual, or confidence\_quant. Those are preserved automatically from the extraction stage. Your JSON output contains ONLY: norm\_subject, norm\_act, condition\_of\_application, norm\_articulation, and role\_rationale.

\#\#\# Contrastive Examples

**Example 1:**
Input:
  norm\_subject: "Mrs. Bennet"
  norm\_act: "promote Jane's extended stay at Netherfield"
  condition: "when a daughter is visiting a household with an eligible bachelor"
  articulation: "Mrs. Bennet is expected to promote Jane's extended stay at Netherfield to encourage Mr. Bingley's attachment."

Output:
  norm\_subject: "a mother of unmarried daughters whose social duty includes securing advantageous matches"
  norm\_act: "promote extended social proximity between a daughter and an eligible suitor"
  condition: "when a daughter has occasion to be in the company of a suitable prospect"
  articulation: "A mother of unmarried daughters is expected to promote extended social proximity between a daughter and an eligible suitor when occasion permits."
  role\_rationale: "Mrs. Bennet occupies the position of a mother with daughters of marriageable age in Regency society, where a mother's primary social duty is securing her daughters' futures through marriage. The norm applies to her qua mother-as-matchmaker, not as a specific individual."

**Example 2:**
Input:
  norm\_subject: "Mr. Darcy"
  norm\_act: "condescend to dance with women of lower social standing"
  condition: "at a public ball where eligible women lack partners"
  articulation: "Mr. Darcy ought to condescend to dance with women at the Meryton assembly."

Output:
  norm\_subject: "a gentleman of high social standing attending a public social gathering"
  norm\_act: "engage socially with those of lower rank, including by dancing with women who lack partners"
  condition: "at a public ball or assembly where his rank creates an expectation of social generosity"
  articulation: "A gentleman of high social standing attending a public ball is expected to engage socially with those of lower rank, including by dancing with women who lack partners."
  role\_rationale: "Darcy's refusal to dance at Meryton is censured precisely because his wealth and standing create a noblesse oblige expectation. The role is not merely 'a wealthy man' but 'a man whose social eminence obliges participation' — the functional capacity (obligation to circulate) flows directly from the structural position (high rank at a public gathering)."

**Example 3 (already partially abstracted — enrich the role):**
Input:
  norm\_subject: "a young woman"
  norm\_act: "not travel alone to visit a gentleman's estate"
  condition: "without a chaperone or familial escort"
  articulation: "A young woman must not travel alone to a gentleman's estate without a chaperone."

Output:
  norm\_subject: "an unmarried young woman of genteel standing who remains under her family's social protection"
  norm\_act: "not travel unaccompanied to a gentleman's estate"
  condition: "without a chaperone or familial escort, particularly when no engagement or formal understanding exists"
  articulation: "An unmarried young woman of genteel standing must not travel unaccompanied to a gentleman's estate without a chaperone, particularly when no engagement or formal understanding exists."
  role\_rationale: "The input already used a partial role ('a young woman') but lacked the relational context (unmarried, under family protection) and the condition that makes the norm operative (no formal understanding). The functional capacity — that her reputation depends on visible propriety — is what makes the chaperonage norm binding."

**Example 4 (totalitarian surveillance — first-name-only character):**
Input:
  norm\_subject: "Winston"
  norm\_act: "show resentment"
  condition: "during the exercise session and in the presence of the telescreen"
  articulation: "Winston must not show resentment during the exercise session and in the presence of the telescreen."

Output:
  norm\_subject: "a citizen living under continuous state surveillance whose every expression is monitored for ideological deviance"
  norm\_act: "show resentment"
  condition: "during mandatory collective activities conducted under surveillance"
  articulation: "A citizen living under continuous state surveillance must not show resentment during mandatory collective activities conducted under surveillance."
  role\_rationale: "Winston is not merely 'a man' — he is a subject of a totalitarian regime where emotional self-regulation is a survival requirement enforced by omnipresent monitoring. The norm applies to him qua surveilled citizen: anyone in his position (monitored, suspected of thoughtcrime) would be bound by the same obligation to suppress visible dissent. 'The telescreen' is a plot-specific device generalized to 'surveillance.'"

**Example 5 (moral dilemma under legal threat — ex-convict identity):**
Input:
  norm\_subject: "Jean Valjean"
  norm\_act: "make his choice between living with a hidden crime or facing the consequences of his past"
  condition: "when facing a moral dilemma between internal and external integrity"
  articulation: "Jean Valjean must make his choice between living with a hidden crime or facing the consequences of his past when he faces a moral dilemma between internal and external integrity."

Output:
  norm\_subject: "a person who has rebuilt a respectable life under a concealed identity after serving a criminal sentence"
  norm\_act: "choose between preserving a concealed but socially productive life or surrendering to legal accountability for past offenses"
  condition: "when the concealment threatens an innocent party's freedom or when continued deception conflicts with one's moral convictions"
  articulation: "A person who has rebuilt a respectable life under a concealed identity after serving a criminal sentence must choose between preserving that life or surrendering to legal accountability when the concealment threatens an innocent party or conflicts with one's moral convictions."
  role\_rationale: "Valjean's dilemma is not personal but structural: it applies to any reformed ex-convict who has assumed a new identity and built social standing. The role is 'a reformed person living under concealment' — the functional capacity (social leadership, moral conscience) and the relational context (legal system that would re-imprison him, innocent person at risk) are what make the norm binding."

**Example 6 (marriage and autonomy — first-name-only female character):**
Input:
  norm\_subject: "Dorothea"
  norm\_act: "consider the financial and social implications of her marriage"
  condition: "before she makes her final decision"
  articulation: "Dorothea ought to consider the financial and social implications of her marriage before she makes her final decision."

Output:
  norm\_subject: "a young woman of independent means whose marriage will substantially alter her legal and financial standing"
  norm\_act: "consider the financial and social implications of her marriage"
  condition: "before making a final decision, particularly when the match involves a significant disparity in age, temperament, or social expectation"
  articulation: "A young woman of independent means whose marriage will substantially alter her legal and financial standing ought to consider the financial and social implications before making her final decision, particularly when the match involves a significant disparity."
  role\_rationale: "Dorothea's position is that of a propertied young woman in a society where marriage is a legal and economic transformation — she will lose control of her estate and daily autonomy. The norm applies to her qua woman-with-property-facing-coverture, not as a specific individual. The relational context (an older scholarly husband, concerned family) is generalized to 'significant disparity.'"

**Example 7 (norm that is already well-abstracted — minimal change):**
Input:
  norm\_subject: "a host"
  norm\_act: "provide for the comfort and entertainment of guests"
  condition: "for the duration of their stay"
  articulation: "A host is obligated to provide for the comfort and entertainment of guests for the duration of their stay."

Output:
  norm\_subject: "a host"
  norm\_act: "provide for the comfort and entertainment of guests"
  condition: "for the duration of their stay"
  articulation: "A host is obligated to provide for the comfort and entertainment of guests for the duration of their stay."
  role\_rationale: "The input norm is already fully abstracted — the subject is a functional social role ('a host'), the act is generalizable, and no character names or plot references are present. No rewrite needed."

\#\#\# Rules

- NEVER invent new norms. You are rewriting, not extracting. Preserve the prescriptive content exactly.
- NEVER change the deontic force (prescriptive\_element, normative\_force). A prohibited norm stays prohibited.
- NEVER introduce character names that weren't in the input. Your output must be name-free.
- If the input norm is already fully abstracted (no names, good role description), output it unchanged and note this in role\_rationale.
- The role\_rationale should explain: (a) what social position the character occupies, (b) what relational context activates the norm, and (c) what functional capacity or duty makes the norm binding. 1-3 sentences.
- When multiple characters are named in a single norm (e.g., names in both subject and condition), abstract ALL of them.
\end{prompttext}

\textbf{User Template:}
\begin{prompttext}
Rewrite the following norm using functional social roles instead of character names. Preserve the prescriptive content exactly — only change the character-specific elements.

\#\#\# Source context

Book summary:
{{book\_summary}}

Text chunk the norm was extracted from:
{{article\_text}}

\#\#\# Norm to rewrite

Input norm:
  prescriptive\_element: {{prescriptive\_element}}
  norm\_subject: {{norm\_subject}}
  norm\_act: {{norm\_act}}
  condition\_of\_application: {{condition\_of\_application}}
  normative\_force: {{normative\_force}}
  norm\_articulation: {{norm\_articulation}}
  context: {{context}}
  quality\_flags: {{quality\_flags}}

The quality\_flags field lists which norm components contain character names or plot-specific references (null means no issues detected). Use this to focus your rewrite: fix the flagged fields, and leave unflagged fields unchanged. If quality\_flags is null, the norm is likely already abstract — output it unchanged unless you spot an issue the automated check missed.

Rewrite using functional social roles. If already fully abstracted, output unchanged.

Output valid JSON:
\end{prompttext}
\end{promptbox}

% Auto-generated from norm-consolidation
\begin{promptbox}{Semantic deduplication and consolidation of extracted norms.}
\label{prompt:norm-consolidation}

\textbf{System Prompt:}
\begin{prompttext}
You are a norm consolidation specialist. Your task is to merge a cluster of semantically similar norms — extracted from prescriptive texts using Joseph Raz's anatomy of norms — into a single **canonical norm** that faithfully represents the shared prescriptive content of the cluster.

\#\#\# What you receive

A JSON array of norm objects, each with:
- **norm\_articulation**: The norm stated as a complete sentence
- **prescriptive\_element**: The deontic "ought" (must, must not, shall, etc.)
- **norm\_subject**: The role or class of persons bearing the obligation
- **norm\_act**: The action prescribed or proscribed
- **condition\_of\_application**: Circumstances under which the norm applies (or null)
- **normative\_force**: obligatory / prohibited / permitted / recommended / discouraged
- **context**: The societal sphere (worship, legal, family, etc.)
- **governs\_info\_flow**: Whether the norm regulates information exchange
- **info\_flow\_note**: CI-vocabulary description if applicable
- **confidence\_quant**: Numeric confidence (0–10)

\#\#\# Your task

1. **Identify the shared prescriptive core** across all member norms. These norms were clustered because they express the same or a very similar obligation/prohibition — but may differ in wording, specificity, or framing.

2. **Produce a canonical merged norm** using Raz's 4 components:
   - **prescriptive\_element**: Choose the deontic term that best captures the shared force. If members use "must" and "shall", pick the most representative.
   - **norm\_subject**: Abstract to the most appropriate role-level. If members say "the faithful", "the believer", "every Muslim" — pick the level that covers all without over-generalising. E.g., "the believer" may subsume "every Muslim" and "the faithful" if they come from different traditions, but prefer the more specific term if all members share the same tradition.
   - **norm\_act**: Merge into a single verb phrase that captures the common action. If one says "bear false witness" and another says "give false testimony", the canonical form might be "give false testimony" or "bear false witness" — pick the most precise.
   - **condition\_of\_application**: Keep the broadest condition that all members satisfy. If one says "in court" and another says "in judicial proceedings", use the broader "in judicial proceedings". If some members are unconditional and others conditional, note this difference in the rationale.

3. **Write a canonical\_articulation**: A single complete sentence expressing the merged norm as a tradition would state it. Be concrete and specific.

4. **Produce an abstraction\_map**: A JSON object showing how specific terms from member norms were abstracted. Keys are the canonical (merged) terms; values are arrays of the original terms they subsume. Only include entries where actual abstraction occurred (i.e., where terms differ across members). Example:
   ```json
   \{
     "norm\_subject": \{"canonical": "the believer", "originals": ["the faithful", "every Muslim", "the believer"]\},
     "norm\_act": \{"canonical": "give false testimony", "originals": ["bear false witness", "give false testimony", "testify falsely"]\}
   \}
   ```

5. **Write a consolidation\_rationale**: 1–3 sentences explaining why these norms were merged and what (if any) nuance was lost in consolidation.

\#\#\# Rules

- **All norms in a cluster share the same normative\_force (deontic operator).** The canonical norm MUST preserve this exact normative\_force — do NOT change it. An "obligatory" cluster stays "obligatory"; a "prohibited" cluster stays "prohibited". The deontic operator is definitional to a norm's identity.
- Do NOT invent norms. The canonical norm must be a faithful generalisation of the cluster members, not a new prescription.
- Do NOT merge norms that are actually distinct obligations. If the cluster contains norms that address genuinely different prescriptions, flag this in the rationale and pick the dominant norm as canonical.
- Prefer specificity over vagueness. Abstract only as much as necessary to subsume all members.
- For context, governs\_information\_flow, and information\_flow\_note: use the majority value. If tied, prefer the value from the highest-confidence member.
- Output valid JSON matching the schema exactly.
\end{prompttext}

\textbf{User Template:}
\begin{prompttext}
The following {{cluster\_size}} norms have been identified as semantically similar and should be consolidated into a single canonical norm.

Member norms:
{{member\_norms}}

Merge these into a single canonical norm following Joseph Raz's anatomy. Produce:
1. The canonical 4-component Raz tuple (prescriptive\_element, norm\_subject, norm\_act, condition\_of\_application)
2. A canonical\_articulation (one complete sentence)
3. An abstraction\_map showing which original terms were abstracted
4. A consolidation\_rationale explaining the merge

Output valid JSON:
\end{prompttext}
\end{promptbox}

\subsection{SFT \& GRPO Training Prompts}

% Auto-generated from grpo-ci-extraction
\begin{promptbox}{CI flow extraction instruction used for SFT and GRPO training.}
\label{prompt:grpo-ci-extraction}

\textbf{System Prompt:}
\begin{prompttext}
You are an expert in the Contextual Integrity (CI) privacy framework, specializing in identifying and classifying **information flows** — instances where information is transmitted, disclosed, concealed, or withheld between agents.
Your task is to analyze the following text passage for information flows using the Contextual Integrity framework.

Before extracting specific information flows, briefly reason about:
1. What social context(s) are present in this passage?
2. What informational norms (rules about appropriate information sharing)
   would a reasonable person expect to apply in these contexts?
3. Are there any information transfers that might violate or conform to
   these contextual norms?

Then identify senders, recipients, subjects, information types, and
transmission principles for each information flow, and provide a
structured extraction.

Respond with ONLY a JSON object in this exact format (no markdown, no extra text):
\{
  "reasoning": "<narrative reasoning trace covering all flows>",
  "has\_information\_exchange": true,
  "flows": [
    \{
      "sender": "...",
      "recipient": "...",
      "subject": "...",
      "information\_type": "...",
      "transmission\_principle": "...",
      "context": "...",
      "appropriateness": "appropriate|inappropriate|ambiguous",
      "norms\_invoked": ["..."],
      "norm\_source": "explicit|implicit|both",
      "is\_new\_flow": false,
      "confidence": 8
    \}
  ]
\}

If there are no information flows, set "has\_information\_exchange" to false and "flows" to [].
\end{prompttext}

\textbf{User Template:}
\begin{prompttext}
{{instruction}}

{{chunk\_text}}
\end{prompttext}
\end{promptbox}

% Auto-generated from grpo-reward-judge
\begin{promptbox}{Normative grounding reward judge for GRPO training.}
\label{prompt:grpo-reward-judge}

\textbf{System Prompt:}
\begin{prompttext}
You are an expert in Helen Nissenbaum's Contextual Integrity framework.
You evaluate whether a specific extracted information flow is grounded in
a set of norms from a normative universe.

You assess two independent criteria:

1. NORM AWARENESS: The model's extraction includes a "norms\_invoked" field
   listing the norms it believes apply to this flow. Do any of those invoked
   norms semantically match the provided norms from the universe?
   Semantic equivalence is sufficient — exact wording match is not required.
   Score from 0.0 (no match at all) to 1.0 (strong semantic match).

2. FLOW GOVERNANCE: Independently of what norms the model invoked, is this
   information flow actually governed by any of the provided norms?
   "Governed" means the norm regulates, constrains, or establishes
   expectations about information flows of this type — between these kinds
   of actors, about this kind of information, in this context.
   Score from 0.0 (flow is unrelated to all norms) to 1.0 (directly governed).

Also assess whether the flow's appropriateness judgment
(appropriate/inappropriate/ambiguous) is consistent with the governing norm.

Respond with a JSON object matching the required schema.
\end{prompttext}

\textbf{User Template:}
\begin{prompttext}
\#\# Source Text
{{chunk\_text}}

\#\# Extracted Information Flow
{{flow\_json}}

\#\# Retrieved Norms (top-k from normative universe)
{{norm\_universe\_json}}

\#\# Evaluation
Assess both criteria independently for this specific information flow:

1. **Norm awareness** (norm\_match\_score): Do the flow's norms\_invoked
   match any of the retrieved norms? Score 0.0–1.0.

2. **Flow governance** (governance\_score): Is this flow governed by any
   of the retrieved norms, regardless of what the model invoked? Score 0.0–1.0.

3. **Appropriateness consistency**: Is the flow's appropriateness judgment
   consistent with the governing norm?

Provide your evaluation as a JSON object.
\end{prompttext}
\end{promptbox}

% Auto-generated from grpo-no-flow-judge
\begin{promptbox}{Coverage judge for no-flow predictions in GRPO reward.}
\label{prompt:grpo-no-flow-judge}

\textbf{System Prompt:}
\begin{prompttext}
You are an expert in Helen Nissenbaum's Contextual Integrity framework.
You assess whether a text passage contains information flows — transfers,
disclosures, or exchanges of information between agents — that fall under
the governance of a given set of norms.

An "information flow" is any instance where information moves from one
agent to another: disclosures, requests, observations, transmissions,
surveillance, data collection, etc.

"Governed" means the norm regulates, constrains, or establishes
expectations about information flows of that type — between those kinds
of actors, about that kind of information, in that context.

Score from 0.0 (passage has no information flows governed by these norms)
to 1.0 (passage clearly contains information flows directly governed by
these norms).

Respond with a JSON object matching the required schema.
\end{prompttext}

\textbf{User Template:}
\begin{prompttext}
\#\# Source Text
{{chunk\_text}}

\#\# Norms from Normative Universe (top-k retrieved)
{{norm\_universe\_json}}

\#\# Assessment
Does this passage describe any information flows (transfers of information
between agents) that are governed by the provided norms?

Consider:
1. Are there any explicit or implicit transfers of information in this passage?
2. Do any of the provided norms regulate the kind of information exchange
   described or implied in the passage?

Provide your assessment as a JSON object with:
- passage\_contains\_governed\_flows (bool)
- coverage\_score (0.0-1.0)
- explanation (brief)
\end{prompttext}
\end{promptbox}

% Auto-generated from grpo-norm-judgment
\begin{promptbox}{Norm application judgment prompt (CIRL vignette evaluation).}
\label{prompt:grpo-norm-judgment}

\textbf{System Prompt:}
\begin{prompttext}
You are an expert in privacy norms and appropriate information sharing.
When presented with a scenario involving the transmission, disclosure,
or withholding of information between people, you assess whether a
described action is appropriate given the social context and the norms
that govern information flows in that setting.
\end{prompttext}

\textbf{User Template:}
\begin{prompttext}
Consider the following situation:

{{scenario}}

Question: Should {{subject}} {{act}} in this context?

Think carefully about what norms of appropriate information sharing
apply in this social context. Consider who the sender and recipient
are, what type of information is involved, and what expectations
govern its transmission.

Respond with ONLY a JSON object in this exact format (no markdown, no extra text):
\{"judgment": "yes" or "no", "reasoning": "your explanation", "norms\_considered": ["norm 1", "norm 2"]\}
\end{prompttext}
\end{promptbox}
